\chapter{INTRODUCTION}

Bridge Weigh-In-Motion (B-WIM) systems provide a non-intrusive method for
estimating the gross vehicle weight (GVW) of heavy goods vehicles (HGVs) as they
cross a bridge at normal traffic speed \cite{Moses1979}. Unlike static weigh
stations, which require vehicles to stop and are costly to operate, B-WIM systems
infer axle loads continuously from the dynamic strain response of an instrumented
bridge structure. This makes them highly practical for traffic enforcement,
infrastructure load monitoring, and pavement design.

A fundamental step in every B-WIM pipeline is \emph{axle detection}: identifying
the number of vehicle axles and, critically, the precise sample positions at
which each axle passes over a reference cross-section of the bridge. Accurate
axle timing is a prerequisite for weight estimation --- an incorrectly detected
or missed axle directly corrupts the subsequent load calculation
\cite{Soberl2025}.

Traditional axle detection methods apply rule-based signal processing: the bridge
strain signal is smoothed (e.g., using a Savitzky-Golay filter
\cite{SavitzkyGolay1964}), and local maxima are then located using a peak-finding
algorithm with manually tuned thresholds. While straightforward to implement,
these approaches depend on hand-crafted parameters and may fail when the
signal-to-noise ratio is low, vehicle speed varies, or multiple closely-spaced
axles produce overlapping strain pulses.

Deep learning models, in particular one-dimensional (1D) convolutional
architectures, have demonstrated strong performance across a wide range of
sequential detection and segmentation tasks. Their ability to learn
task-specific representations directly from raw data, without explicit feature
engineering, makes them attractive candidates for the axle detection problem.

\section{OBJECTIVES}

The main objectives of this bachelor thesis are:

\begin{enumerate}
    \item Implement and evaluate a classical signal-processing baseline for axle
          detection using Savitzky-Golay smoothing and peak detection.
    \item Design, train, and evaluate a 1D U-Net Convolutional Neural Network
          (CNN) for per-sample binary axle detection.
    \item Design, train, and evaluate a Temporal Convolutional Network (TCN) with
          dilated residual blocks, ensuring the receptive field covers the full
          signal length.
    \item Compare all three approaches on the same held-out test set using
          consistent, axle-level evaluation metrics.
    \item Identify the limitations of the current study and propose directions
          for future work.
\end{enumerate}

\section{RESEARCH QUESTIONS}

\begin{enumerate}
    \item Can deep learning models accurately detect axle positions from B-WIM
          strain signals despite severe class imbalance (287:1 ratio)?
    \item How does a deep learning approach compare to a classical peak-detection
          baseline in terms of F1-score and timing accuracy?
    \item Does a lightweight TCN outperform a larger U-Net CNN on this task?
\end{enumerate}

\section{THESIS STRUCTURE}

Chapter~2 provides background on B-WIM systems and reviews related work on axle
detection and deep learning for time series. Chapter~3 describes the dataset, all
three methods, the training procedure, and the evaluation protocol. Chapter~4
presents and discusses the experimental results. Chapter~5 concludes the thesis
and outlines directions for future work.
